\documentclass[11pt]{article}

\usepackage[margin=1.8cm]{geometry}
\usepackage[T1]{fontenc}
\usepackage{lmodern}
\usepackage{xcolor}
\usepackage{chemfig}
\usepackage{chemformula}
\usepackage{subcaption}
\usepackage{caption}

\pagestyle{empty}
\captionsetup{font=small,labelfont=bf}
\captionsetup[subfigure]{labelformat=parens,font=small,labelfont=bf}

\definecolor{backbonegray}{RGB}{120,135,135}

\setchemfig{
  atom sep=1.95em,
  bond offset=2pt,
  bond style={line width=0.85pt},
  cram width=3.5pt,
  cram dash width=0.45pt,
  cram dash sep=1.1pt,
  double bond sep=2.2pt
}

\newcommand{\moltitle}[1]{\vspace{0.35em}\textbf{#1}}
\newcommand{\bbatom}[1]{\textcolor{backbonegray}{#1}}
\newlength{\peroxidepanelheight}
\setlength{\peroxidepanelheight}{12.7em}

\begin{document}

\begin{figure}
\centering

\begin{subfigure}[t]{0.46\textwidth}
\centering
\schemestart
\chemfig{
  @{cysN}\bbatom{N}
    (-[:120,,,,draw=backbonegray] @{cysHN1}\bbatom{H})
    (-[:240,,,,draw=backbonegray] @{cysHN2}\phantom{H})
  -[,,,,draw=backbonegray] @{cysCA}\bbatom{CA}
    (
      -[2,,,,draw=backbonegray] @{cysH}\bbatom{HA}
    )
    (
      -[6,1.35] @{cysCB}CB
        (-[4,1.25] @{cysHB1}HB1)
        (-[0,1.25] @{cysHB2}HB2)
        -[6,1.35] @{cysSG}{SG}
          -[6,1.25] @{cysHG}HG1
    )
  -[,,,,draw=backbonegray] @{cysC}\bbatom{C}
    (=[1,,,,draw=backbonegray] @{cysO}\bbatom{O})
    -[-1,,,,draw=backbonegray] @{cysOH}\phantom{O}
}
\schemestop
\caption{Cysteine}
\end{subfigure}
\hfill
\begin{subfigure}[t]{0.46\textwidth}
\centering
\schemestart
\chemfig{
  @{secN}\bbatom{N}
    (-[:120,,,,draw=backbonegray] @{secHN1}\bbatom{H})
    (-[:240,,,,draw=backbonegray] @{secHN2}\phantom{H})
  -[,,,,draw=backbonegray] @{secCA}\bbatom{CA}
    (
      -[2,,,,draw=backbonegray] @{secH}\bbatom{HA}
    )
    (
      -[6,1.35] @{secCB}CB
        (-[4,1.25] @{secHB1}HB1)
        (-[0,1.25] @{secHB2}HB2)
        -[6,1.35] @{secSE}{SE}
          -[6,1.25] @{secHG}HG1
    )
  -[,,,,draw=backbonegray] @{secC}\bbatom{C}
    (=[1,,,,draw=backbonegray] @{secO}\bbatom{O})
    -[-1,,,,draw=backbonegray] @{secOH}\phantom{O}
}
\schemestop
\caption{Selenocysteine}
\end{subfigure}

\vspace{1.6em}

\begin{subfigure}[t]{0.46\textwidth}
\centering
\begin{minipage}[c][\peroxidepanelheight][c]{\linewidth}
\centering
  \schemestart
  \chemfig{
    @{hpa}H1
    - @{hpoa}{O1}
    - @{hpob}{O2}
    - @{hpb}H2
  }
  \schemestop
\end{minipage}
\caption{Hydrogen peroxide}
\end{subfigure}
\hfill
\begin{subfigure}[t]{0.46\textwidth}
\centering
\schemestart
\chemfig{
  @{glnN}\bbatom{N}
    (-[:120,,,,draw=backbonegray] @{glnHN1}\bbatom{H})
    (-[:240,,,,draw=backbonegray] @{glnHN2}\phantom{H})
  -[,,,,draw=backbonegray] @{glnCA}\bbatom{CA}
    (
      -[2,,,,draw=backbonegray] @{glnH}\bbatom{HA}
    )
    (
      -[6,1.28] @{glnCB}CB
        (-[4,1.18] @{glnHB1}HB1)
        (-[0,1.18] @{glnHB2}HB2)
        -[6,1.28] @{glnCG}CG
          (-[4,1.18] @{glnHG1}HG1)
          (-[0,1.18] @{glnHG2}HG2)
          -[6,1.28] @{glnCD}CD
            (=[0,1.18] @{glnOE}OE1)
            -[6,1.18] @{glnNE}NE2
              (-[4,1.12] @{glnHE21}HE21)
              (-[0,1.12] @{glnHE22}HE22)
    )
  -[,,,,draw=backbonegray] @{glnC}\bbatom{C}
    (=[1,,,,draw=backbonegray] @{glnO}\bbatom{O})
    -[-1,,,,draw=backbonegray] @{glnOH}\phantom{O}
}
\schemestop
\caption{Glutamine}
\end{subfigure}

\end{figure}

\end{document}
